\section{Introduction}
\label{sec:introduction}

Tax-benefit microsimulation models apply policy rules to household microdata to estimate fiscal and
distributional effects. At the national level, that usually means starting from a survey such as the
Current Population Survey (CPS), improving measured incomes and program participation, and adjusting
weights so the dataset reproduces known aggregates. Subnational work is harder. Analysts often want
estimates for states, congressional districts, or local areas, but the underlying survey was not
designed to support all of those geographies at once.

That gap matters in practice, because many reforms are proposed, administered, or debated below the
national level. A state tax proposal needs state-specific incomes, filing patterns, and benefit
participation. District-level analysis needs those same outcomes to aggregate correctly within each
district while staying consistent with state and national totals. A usable subnational dataset
therefore has to satisfy a hierarchical calibration problem rather than a single national one.

One way to meet that demand is to build a dataset that carries as much of the relevant variation as
possible. A pipeline can combine several surveys and administrative files, impute the variables that
any single source lacks, and place records across geographies so that local detail is present in the
candidate data. Each of these steps adds records, or variants of records. The result is a candidate
dataset rich enough to represent the variation the targets describe, and large enough that it
becomes expensive to store, calibrate, and simulate. A file with millions of active records supports
detailed subnational work but is slow to run; a much smaller file is easier to ship, but only if the
loss in fidelity is acceptable.

The candidate dataset therefore has to be reduced to a size that can be shipped and simulated, and
this paper treats that reduction as a sampling problem: given the large candidate universe and a
fixed record budget, which records should the final dataset keep, and what weight should each kept
record carry? The simplest answer is to draw a random subset
and reweight it to the targets. We study an alternative in which the selection itself is informed by
the targets.

The method is $L_0$ regularization with the Hard Concrete relaxation \citep{louizos2018}. Each
candidate record receives a gate that the optimizer can drive toward zero or one, and the expected
number of open gates enters the objective directly. Minimizing calibration error and the gate count
together selects a subset of records and fits their weights in one optimization. The retained-record
count becomes a quantity the optimizer controls rather than a fixed draw, and a separate weight
penalty limits how concentrated the fitted weights can become. In practical terms, one framework can
target a compact national dataset or a larger subnational dataset, at a chosen record budget, with
weights kept inside a usable range.

We situate this work in \populace{} \citep{populace2026}, PolicyEngine's pipeline for microsimulation
dataset construction. \populace{} passes a single weighted data structure through a sequence of
modular stages---source loading, imputation, geographic assignment, target construction, and
calibration---each implemented as a separate package operating on that shared structure. Because the
stages are modular, the data sources, the imputation model, the geographic assignment, and the
calibration method are interchangeable components rather than fixed code, and the calibration step
studied here is one such component. The paper describes the pipeline at the level needed to follow the
calibration step, and then describes the configuration used here.

This paper is a proof of concept on \populace{}'s current target surface: it asks whether
$L_0$ selection with Hard Concrete gates can produce a sparse, calibrated dataset on the targets
\populace{} assembles today. Local-area and subnational calibration are the production motivation for
the method and a future application of it, rather than results this paper presents.

The empirical question is whether informed selection earns its added complexity. We ask three
things. First, does $L_0$ selection reproduce the targets more accurately than random sampling,
survey-weight sampling, and combinatorial optimisation at the same record budget? Second, does the retained sample generalize, in the sense that it also
matches targets held out of the fit? Third, does the framework give usable control over the size of
the output, the geographic level it focuses on, and the magnitude of the fitted weights?

The rest of the paper follows that structure. Section~\ref{sec:background} places the work in the
literature on spatial microsimulation, survey calibration, and sparse selection.
Section~\ref{sec:data} describes the base microdata, the \populace{} pipeline, and the
administrative targets. Section~\ref{sec:methodology} presents the calibration method and the
sampling experiment design. Section~\ref{sec:results} reports the comparison against the sampling
baselines and the behaviour of the size, geography, and weight controls. Section~\ref{sec:discussion}
interprets those results and discusses their limits, and Section~\ref{sec:conclusion} concludes.
